\section{The Weierstrass Theorem}

Recall that in definition~\ref{def:funzione_limitata} we defined the concepts of
maximum, minimum, supremum, and infimum of a real-valued function.

\begin{lemma}[minimizing/maximizing sequences]%
\mymargin{minimizing sequences}%
\index{sequence!minimizing}%
\index{sequence!maximizing}%
Let $A$ be a non-empty set and let $f\colon A \to \RR$ be a function. Then there
exist two sequences $a_n$ and $b_n$ of points in $A$ such that
\[
  \lim_{n\to +\infty} f(a_n) = \inf f(A), \qquad
  \lim_{n\to +\infty} f(b_n) = \sup f(A).
\]
\end{lemma}
\mymargin{minimizing sequences}
%
\begin{proof}
Recall that $f(A) = \ENCLOSE{f(x)\colon x \in A}$ is the image of the function
$f$. We will prove the statement for the infimum; an analogous result can be
obtained for the supremum.

Let $m=\inf f(A)$.
If $m=-\infty$, it means that $f(A)$ is not bounded below, in particular for
every $n\in \RR$ there exists $a_n$ such that $f(a_n) < - n$.
Thus (by comparison) $f(a_n) \to -\infty$ as we wanted to demonstrate.

If $m\in \RR$, by the characterizing properties of the infimum, we know that for
every $\eps>0$ there exists $a\in A$ such that $f(a) < m + \eps$.
For every $n\in\NN$, we can choose $\eps=1/n$ and thus obtain a sequence $a_n$
such that $f(a_n) < m + 1/n$.
On the other hand, since $m$ is a lower bound of $f(A)$, we know that $m \le
f(a_n)$.
We thus have $m \le f(a_n) < m+ 1/n$, and by the squeeze theorem, we can
conclude that $f(a_n) \to m$ as $n\to +\infty$.
\end{proof}

\begin{theorem}[Weierstrass]%
\label{th:weierstrass}%
\mymark{***}%
\mymargin{Weierstrass theorem}%
\index{theorem!Weierstrass}%
\index{Weierstrass!theorem}%
\mynote{see historical notes on page~\pageref{note:isoperimetrico}}%
Let $a,b\in \RR$, $a\le b$, and $f\colon [a,b]\to \RR$ be a continuous function.
Then there exist points of maximum and minimum for $f$ on $[a,b]$.
\end{theorem}
%
\begin{proof}
\mymark{***}%
We will only prove that $f$ has a minimum; the proof for the maximum proceeds in
an entirely analogous manner.

Let $m= \inf f([a,b])$.
By the previous lemma, we know that there exists a minimizing sequence $a_n$,
i.e., such that $a_n \in [a,b]$ and $f(a_n)\to m$ as $n\to +\infty$.

By the Bolzano-Weierstrass theorem, from the sequence $a_n$, we can extract a
convergent subsequence $a_{n_k}$: $a_{n_k} \to x_0$.
Since $a_{n_k} \in [a,b]$, by the theorem of sign permanence, we also have $x_0
\in [a,b]$ (apply the sign permanence to the sequences $a_{n_k}-a$ and
$b-a_{n_k}$).

Thus, we have a sequence $a_{n_k}\to x_0$ with $a_{n_k}\in [a,b]$ and $x_0 \in
[a,b]$. Since $f$ is continuous, we will have $f(a_{n_k}) \to f(x_0)$.
But we knew that $f(a_n)\to m$, and thus also $f(a_{n_k}) \to m$.
We conclude that $f(x_0) = m$, i.e., $m$, the infimum, is a value assumed by the
function and is therefore a minimum.
On its part, $x_0$ is a point of absolute minimum.
\end{proof}

\begin{corollary}[boundedness of continuous functions]
Let $f\colon [a,b]\to \RR$ be a continuous function. Then $f$ is bounded.
\end{corollary}
\begin{proof}
Since $f$ has a maximum $M$ and a minimum $m$, we have $f(x)\in [m,M]$ for every
$x\in[a,b]$.
Obviously, $m>-\infty$ and $M<+\infty$ since $m$ and $M$ are values of the
function $f$.
\end{proof}

It often happens that one needs to determine whether a continuous function
$f\colon I\to \RR$ defined on an interval $I$ assumes maximum or minimum values,
but usually, the interval $I$ is not closed and bounded. More frequently, one
falls into this generalized form of the Weierstrass theorem.
\begin{theorem}[Generalized Weierstrass]%
  \label{th:weierstrass_generalizzato}%
  \mymark{***}%
  \mymargin{generalized Weierstrass theorem}%
  \index{theorem!generalized Weierstrass}%
  \index{Weierstrass!generalized theorem}%
Let $a\ge -\infty$, $b\le +\infty$, with $a<b$, and let $f\colon (a,b)\to \RR$
be a continuous function.
If 
\[
 \lim_{x\to a^+} f(x) = +\infty, \qquad 
 \lim_{x\to b^+} f(x) = +\infty
\]
then $f$ assumes a minimum value at at least one point in $(a,b)$.
\end{theorem}
\begin{proof}
Let $m=\inf f((a,b))$.
By the previous lemma, there exists a minimizing sequence $a_n$, i.e., such that
$a_n \in (a,b)$ and $f(a_n)\to m$ as $n\to +\infty$.
By the Bolzano-Weierstrass theorem~\ref{th:bolzano_weierstrass}, from the
sequence $a_n$, we can extract a convergent subsequence $a_{n_k}$: $a_{n_k} \to
x_0$.
By the sign permanence theorem~\ref{th:permanenza_segno}, it must be $x_0 \in
[a,b]$.
But if $x_0=a$, we would have $f(a_{n_k}) \to +\infty$, contradicting the fact
that $a_n$ is a minimizing sequence.
The same would happen if $x_0=b$. Thus $x_0 \in (a,b)$.
We conclude as in the Weierstrass theorem, by the continuity of $f$ at $x_0$, we
deduce $f(x_0) = m$, and thus $x_0$ is a point of minimum.
\end{proof}

\begin{exercise}
In each of the following cases, decide whether there exists a function with the
indicated characteristics:
\begin{enumerate}
  \item $f\colon \RR\to (0,1)$ continuous and bijective;
  \item $f\colon \RR\to [0,1]$ continuous and surjective;
  \item $f\colon [0,1]\to \RR$ continuous and surjective;
  \item $f\colon (0,1]\to \RR$ continuous and surjective;
  \item $f\colon (0,1]\to \RR$ bijective. 
\end{enumerate}
\end{exercise}
%%%%%%%%%%%
%%%%%%%%%%%